\documentclass[10pt,a4paper]{amsart}
\usepackage{amsmath,amssymb,amsthm,graphicx}
% Retained minimal preamble for the Stacks Project chapter sites-cohomology.tex (Cohomology on Sites).
% Verbatim excerpt of the chapter preamble: Xy-pic setup (native preamble lines 14-18) and the native
% theorem declarations with the mathematical macros (native preamble lines 154-217). Copied mechanically;
% nothing added, removed or redefined. The wrapper class and the amsmath/amssymb/amsthm/graphicx base
% packages are supplied by the assembler, so the class/externaldocument part of the native preamble is omitted.
% For commutative diagrams we use Xy-pic
\usepackage[all]{xy}

% We use 2cell for 2-commutative diagrams.
\xyoption{2cell}
% Theorem environments.
%
\theoremstyle{plain}
\newtheorem{theorem}[subsection]{Theorem}
\newtheorem{proposition}[subsection]{Proposition}
\newtheorem{lemma}[subsection]{Lemma}

\theoremstyle{definition}
\newtheorem{definition}[subsection]{Definition}
\newtheorem{example}[subsection]{Example}
\newtheorem{exercise}[subsection]{Exercise}
\newtheorem{situation}[subsection]{Situation}

\theoremstyle{remark}
\newtheorem{remark}[subsection]{Remark}
\newtheorem{remarks}[subsection]{Remarks}

\numberwithin{equation}{subsection}

% Macros
%
\def\lim{\mathop{\mathrm{lim}}\nolimits}
\def\colim{\mathop{\mathrm{colim}}\nolimits}
\def\Spec{\mathop{\mathrm{Spec}}}
\def\Hom{\mathop{\mathrm{Hom}}\nolimits}
\def\Ext{\mathop{\mathrm{Ext}}\nolimits}
\def\SheafHom{\mathop{\mathcal{H}\!\mathit{om}}\nolimits}
\def\SheafExt{\mathop{\mathcal{E}\!\mathit{xt}}\nolimits}
\def\Sch{\mathit{Sch}}
\def\Mor{\mathop{\mathrm{Mor}}\nolimits}
\def\Ob{\mathop{\mathrm{Ob}}\nolimits}
\def\Sh{\mathop{\mathit{Sh}}\nolimits}
\def\NL{\mathop{N\!L}\nolimits}
\def\CH{\mathop{\mathrm{CH}}\nolimits}
\def\proetale{{pro\text{-}\acute{e}tale}}
\def\etale{{\acute{e}tale}}
\def\QCoh{\mathit{QCoh}}
\def\Ker{\mathop{\mathrm{Ker}}}
\def\Im{\mathop{\mathrm{Im}}}
\def\Coker{\mathop{\mathrm{Coker}}}
\def\Coim{\mathop{\mathrm{Coim}}}

% Boxtimes
%
\DeclareMathSymbol{\boxtimes}{\mathbin}{AMSa}{"02}

%
% Macros for moduli stacks/spaces
%
\def\QCohstack{\mathcal{QC}\!\mathit{oh}}
\def\Cohstack{\mathcal{C}\!\mathit{oh}}
\def\Spacesstack{\mathcal{S}\!\mathit{paces}}
\def\Quotfunctor{\mathrm{Quot}}
\def\Hilbfunctor{\mathrm{Hilb}}
\def\Curvesstack{\mathcal{C}\!\mathit{urves}}
\def\Polarizedstack{\mathcal{P}\!\mathit{olarized}}
\def\Complexesstack{\mathcal{C}\!\mathit{omplexes}}
% \Pic is the operator that assigns to X its picard group, usage \Pic(X)
% \Picardstack_{X/B} denotes the Picard stack of X over B
% \Picardfunctor_{X/B} denotes the Picard functor of X over B
\def\Pic{\mathop{\mathrm{Pic}}\nolimits}
\def\Picardstack{\mathcal{P}\!\mathit{ic}}
\def\Picardfunctor{\mathrm{Pic}}
\def\Deformationcategory{\mathcal{D}\!\mathit{ef}}
\usepackage{cleveref}
\crefname{originallemma}{original lemma}{original lemmas}
\Crefname{originallemma}{Original lemma}{Original lemmas}
\makeatletter
\newlabel{more-algebra-lemma-open-mapping}{{0CQW}{source}{}{sourcepointer-1}{}}
\newlabel{more-algebra-lemma-open-mapping@cref}{{[originallemma][0][]0CQW}{[1][0][]source}{}{}{}}
\makeatother
\title{The limit topology on the cohomology of an inverse system of sheaves is I-adic}
\author{The Stacks Project Authors}
\date{stacks-project:sites-cohomology.tex:local-snapshot-20261011; GNU FDL 1.2 or later}
\begin{document}
\maketitle
\noindent\textit{Editorial note.} All delivered mathematical text is the author's own, copied from the pinned Stacks Project chapter file sites-cohomology.tex; no statement or proof step was rewritten, reordered or removed.\par\medskip
\noindent\textit{Editorial note.} The author marks one structural verification as omitted: at source line 4905 the sentence recording that the direct sum of the E\_n is a graded submodule and that the map to the direct sum of the F\textasciicircum{}n/F\textasciicircum{}\{n+1\} is graded ends with 'details omitted'. That verification concerns the bookkeeping of the graded structures, and the load-bearing steps of the argument, the finite generation, the identity IF\textasciicircum{}n = F\textasciicircum{}\{n+1\} and the open mapping lemma, are all written out inside the counted span. The omission is recorded here rather than repaired.\par\medskip
\noindent\textit{Editorial note.} The counted statement carries the author's footnote on the uniform form of the conclusion; it lies inside the statement span and is delivered as written.\par\medskip
\noindent\textit{Editorial note.} Cross-chapter restatement: the same label occurs in chapter cohomology for a topological space. That statement is the special case of this one, so per the lane-s4 rule only this more general occurrence is counted here.\par\medskip
\noindent\textit{Editorial note.} The single labelled reference occurring inside the counted spans points to another chapter of the Project; it is carried in assembly.external\_references as a non-load-bearing navigation pointer recorded at the invoking line.\par\medskip
\noindent\textit{Editorial note.} The preamble is the chapter's own retained minimal preamble (sources/stacks-sites-cohomology/retained-preamble.tex), a verbatim excerpt of the native preamble: the Xy-pic setup and the native theorem declarations with the mathematical macros.\par\medskip
\noindent\textit{Allowed standard prerequisites (editorial).}\begin{itemize}
\item I-adic and limit topologies on the limit of an inverse system of modules. Same chapter, the comparison of the I-adic topology and the limit topology on M = lim H\textasciicircum{}p(C, F\_n). Supplies both topologies whose equality is asserted, and the chain F\textasciicircum{}n = ker(M -> H\textasciicircum{}p(C, F\_n)) of neighbourhoods of the limit topology.
\item Graded module structure over the ring of the powers of I. Same chapter, the graded module structure of the direct sum of the H\textasciicircum{}p(C, I\textasciicircum{}n F\_\{n+1\}) over the direct sum of the I\textasciicircum{}n/I\textasciicircum{}\{n+1\}. Supplies the ascending chain condition hypothesis and the graded generators used in the proof.
\item Open mapping lemma for complete linearly topologized modules. More on Algebra, the open mapping lemma applied to u : (F\textasciicircum{}n)\textasciicircum{}\{direct sum r\} -> F\textasciicircum{}\{n+1\}. Supplies the step that upgrades density of the image to surjectivity of u.
\end{itemize}\medskip
\ifdefined\hypertarget\hypertarget{sourcepointer-1}{}\fi
\noindent\textit{Source reference.} Original lemma 0CQW: sources/stacks-sites-cohomology/pinned-original-source.tex, source lines 4935--4935. Cross-chapter navigation pointer invoked at proof lines 4934-4935, where the author applies the open mapping lemma for complete linearly topologized modules to conclude that the map u on the dense image is open. The target chapter is not part of this card's source tree, so the pointer records the invoking line; it only resolves the reference for the reader, and no bibliography citation is involved.\par\medskip
\begin{lemma}
\label{lemma-topology-I-adic-general}
Let $I$ be an ideal of a ring $A$. Let $\mathcal{C}$ be a site. Let
$$
\ldots \to \mathcal{F}_3 \to \mathcal{F}_2 \to \mathcal{F}_1
$$
be an inverse system of sheaves of $A$-modules on $\mathcal{C}$ such that
$\mathcal{F}_n = \mathcal{F}_{n + 1}/I^n\mathcal{F}_{n + 1}$.
Let $p \geq 0$. Assume
$$
\bigoplus\nolimits_{n \geq 0} H^p(\mathcal{C}, I^n\mathcal{F}_{n + 1})
$$
satisfies the ascending chain condition as a graded
$\bigoplus_{n \geq 0} I^n/I^{n + 1}$-module.
Then the limit topology on $M = \lim H^p(\mathcal{C}, \mathcal{F}_n)$
is the $I$-adic topology.
\end{lemma}

\begin{proof}
Set $F^n = \Ker(M \to H^p(\mathcal{C}, \mathcal{F}_n))$ for $n \geq 1$ and
$F^0 = M$.
Observe that $I F^n \subset F^{n + 1}$. In particular $I^n M \subset F^n$.
Hence the $I$-adic topology is finer than the limit topology. For
the converse, we will show that given $n$
there exists an $m \geq n$ such that $F^m \subset I^nM$\footnote{In fact,
there exist a $c \geq 0$ such that $F^{n + c} \subset I^nM$ for all $n$.}.
We have injective maps
$$
F^n/F^{n + 1} \longrightarrow H^p(\mathcal{C}, \mathcal{F}_{n + 1})
$$
whose image is contained in the image of
$H^p(\mathcal{C}, I^n\mathcal{F}_{n + 1}) \to
H^p(\mathcal{C}, \mathcal{F}_{n + 1})$.
Denote
$$
E_n \subset H^p(\mathcal{C}, I^n\mathcal{F}_{n + 1})
$$
the inverse image of $F^n/F^{n + 1}$. Then $\bigoplus E_n$ is
a graded $\bigoplus I^n/I^{n + 1}$-submodule of
$\bigoplus H^p(\mathcal{C}, I^n\mathcal{F}_{n + 1})$ and
$\bigoplus E_n \to \bigoplus F^n/F^{n + 1}$ is a homomorphism of graded
modules; details omitted. By assumption $\bigoplus E_n$ is generated by
finitely many homogeneous elements over $\bigoplus I^n/I^{n + 1}$.
Since $E_n \to F^n/F^{n + 1}$ is surjective, we see that
the same thing is true of $\bigoplus F^n/F^{n + 1}$.
Hence we can find $r$ and $c_1, \ldots, c_r \geq 0$ and
$a_i \in F^{c_i}$ whose images in $\bigoplus F^n/F^{n + 1}$ generate.
Set $c = \max(c_i)$.

\medskip\noindent
For $n \geq c$ we claim that $I F^n = F^{n + 1}$. The claim shows that
$F^{n + c} = I^nF^c \subset I^nM$ as desired. To prove the claim
suppose $a \in F^{n + 1}$. The image of
$a$ in $F^{n + 1}/F^{n + 2}$ is a linear combination
of our $a_i$. Therefore $a  - \sum f_i a_i \in F^{n + 2}$
for some $f_i \in I^{n + 1 - c_i}$. Since
$I^{n + 1 - c_i} = I \cdot I^{n - c_i}$ as $n \geq c_i$ we can write
$f_i = \sum g_{i, j} h_{i, j}$ with $g_{i, j} \in I$
and $h_{i, j}a_i \in F^n$. Thus we see that
$F^{n + 1} = F^{n + 2} + IF^n$.
A simple induction argument gives $F^{n + 1} = F^{n + e} + IF^n$
for all $e > 0$. It follows that $IF^n$ is dense in $F^{n + 1}$.
Choose generators $k_1, \ldots, k_r$ of $I$ and consider
the continuous map
$$
u : (F^n)^{\oplus r} \longrightarrow F^{n + 1},\quad
(x_1, \ldots, x_r) \mapsto \sum k_i x_i
$$
(in the limit topology).
By the above the image of $(F^m)^{\oplus r}$ under $u$ is dense in
$F^{m + 1}$ for all $m \geq n$. By the open mapping lemma
(More on Algebra, Lemma \ref{more-algebra-lemma-open-mapping}) we find
that $u$ is open. Hence $u$ is surjective. Hence $IF^n = F^{n + 1}$
for $n \geq c$. This concludes the proof.
\end{proof}

\end{document}
